% !PRINTED 57
\weberpage{57}
%<<CONT

stens nicht, ohne dass noch verschiedene Potenzen von $t$ in den einzelnen Gliedern bleiben.

Durch Vermehrung der Veränderlichen kann man jede nicht homogene Function in eine homogene von gleichem Grade verwandeln. Ist nämlich $m-1$ die Anzahl der Variablen in einer nicht homogenen Function $n$ten Grades, so setzen wir
\[
 x={x_1\over x_m},\qquad y={x_2\over x_m},\qquad z={x_3\over x_m},\ldots,
\]
und erhalten in
\[
 x_m^n F\!\left({x_1\over x_m},{x_2\over x_m},{x_3\over x_m},\ldots\right)
\]
eine ganze homogene Function $n$ten Grades der Variablen $x_1,x_2,\ldots,x_m$, die wir mit
\[
 \Phi(x_1,x_2,\ldots,x_m)
\]
bezeichnen.

Es empfiehlt sich bisweilen, die homogenen Functionen mehrerer Variablen \gesperrt{mit den Polynomialcoefficienten} zu schreiben. Wir setzen daher
\begin{equation}
 \Phi(x_1,x_2,\ldots,x_m)=
 \sum{\Pi(n)\over\Pi(\alpha_1)\Pi(\alpha_2)\cdots\Pi(\alpha_m)}
 A_{\alpha_1,\alpha_2,\ldots,\alpha_m}
 x_1^{\alpha_1}x_2^{\alpha_2}\cdots x_m^{\alpha_m},
\tag{2}
\end{equation}
wo sich die Summe auf alle nicht negativen, der Bedingung
\begin{equation}
 \alpha_1+\alpha_2+\cdots+\alpha_m=n
\tag{3}
\end{equation}
genügenden Zahlen erstreckt. Diese Bezeichnungsweise, ohne die Beschränkung (3), ist auch auf nicht homogene Functionen anwendbar.

Man kann aber die homogene Function auch so darstellen:
\begin{equation}
 \Phi(x_1,x_2,\ldots,x_m)=
 \sum A_{\nu_1,\nu_2,\ldots,\nu_n}x_{\nu_1}x_{\nu_2}\cdots x_{\nu_n},
\tag{4}
\end{equation}
worin jeder der Indices $\nu_1,\nu_2,\ldots,\nu_n$ von den übrigen unabhängig die Werthreihe $1,2,\ldots,m$ zu durchlaufen hat. Die Summe (4) besteht also aus $m^n$ Gliedern, die aber nicht alle von einander verschieden sind. Das Product $x_{\nu_1}x_{\nu_2}\cdots x_{\nu_n}$ bleibt nämlich ungeändert, wenn die Indices $\nu_1,\nu_2,\ldots,\nu_n$ beliebig unter einander permutirt werden. Die Anzahl der Permutationen von $n$ Elementen beträgt aber $\Pi(n)$. Sind unter diesen Elementen je $\alpha_1,\alpha_2,\ldots$ einander gleich, so reducirt sich die Zahl der Permutationen auf
\[
 {\Pi(n)\over \Pi(\alpha_1)\Pi(\alpha_2)\cdots},
\]
%CONT>>
